\section{Features, directions and superposition}
\label{sec:features}

The methods of \cref{sec:observation} work with components that the architecture hands us: heads, layers,
positions. To say what a model \emph{represents} we need a second kind of unit, one that corresponds to a
concept. This section explains why single neurons are not that unit, and what the current working hypothesis
is.

\subsection{Features and the trouble with neurons}

We use the word \emph{feature} for a property of the input, or of the model's intermediate reasoning, that the
model represents internally because it is useful for prediction. Examples are ``the current token is a first
name'', ``this text is French'', ``the subject of the sentence is plural'' or ``this name has occurred before''.
The definition is informal on purpose. Making it precise is an open problem, and the hypotheses below are
attempts to do so.

The simplest conceivable design would give each feature its own neuron. Some neurons do behave like this. Many
do not. \citet{bricken2023monosemanticity} describe a neuron of a small language model that is active on
academic citations, on English dialogue, on HTTP requests and on Korean text. Such neurons are called
\emph{polysemantic}. Reading
them one at a time gives descriptions that are both vague and incomplete, and \citet{bolukbasi2021interpretability}
showed that the handful of strongest activations of a neuron can suggest a clean meaning that disappears when
one looks at a different data set.

For the residual stream the hope was never reasonable. By the last fact of \cref{sec:facts}, its coordinates
have no special role, so a feature has no reason to align with one of them.

\subsection{The linear representation hypothesis}
\label{sec:lrh}

If features are not coordinates, the next simplest possibility is that they are \emph{directions}.

\begin{definition}[Linear representation hypothesis, working form]\label{def:lrh}
For an activation $x\in\R^d$ at a fixed place in the model (say, the stream after a given layer) there are unit
vectors $v_1,\dots,v_N\in\R^d$, called \emph{feature directions}, and non-negative functions $f_1,\dots,f_N$ of
the input, called \emph{feature activations}, such that
\begin{equation}\label{eq:lrh}
x \;\approx\; \sum_{j=1}^{N}f_j\,v_j
\end{equation}
on all inputs of interest, where each $f_j$ measures the presence of one feature and only few of the $f_j$ are
nonzero on any single input.
\end{definition}

The hypothesis says two things: the \emph{presence} of a feature is a direction, and its \emph{strength} is the
length along that direction. It is a hypothesis about trained networks and not a theorem. The evidence for it
comes from several independent sources.

\begin{itemize}
\item \textbf{Vector arithmetic.} In word embeddings, differences between vectors encode relations. The famous
example is that the vector of \tok{king} minus that of \tok{man} plus that of \tok{woman} lies close to the
vector of \tok{queen} \citep{mikolov2013linguistic}.
\item \textbf{Linear probes.} A wide range of properties is linearly decodable (\cref{sec:probing}), including
ones as abstract as the truth of a statement \citep{marks2024geometry}. The concept activation vectors of
\citet{kim2018interpretability} apply the same idea to image classifiers.
\item \textbf{Steering.} Adding a suitable vector to the stream switches a behaviour on or off
(\cref{sec:steering}). This is causal evidence, which the first two items are not.
\item \textbf{An architectural reason.} Every component reads the stream through a linear map. A property that
is stored along a direction can be read by one row of such a map, so linear storage is the cheapest format for
whatever later layers need.
\end{itemize}

\paragraph{Estimating a direction from labelled data.} Suppose we have activations $x$ with labels
$y\in\{0,1\}$ that say whether a concept is present. The \emph{difference of means} is the vector
$\E[x\mid y=1]-\E[x\mid y=0]$, estimated by the two sample averages. Under \eqref{eq:lrh} taken as an equality
we can say exactly what it measures. By linearity of the expectation,
\begin{equation}\label{eq:diffmeans}
\E[x\mid y=1]-\E[x\mid y=0] \;=\; \sum_{j=1}^{N}\big(\E[f_j\mid y=1]-\E[f_j\mid y=0]\big)\,v_j .
\end{equation}
If the label marks the presence of feature $j$ and all other features have the same average in both classes,
only one term survives, and the difference of means is a positive multiple of $v_j$. If some other feature is
more frequent in one class, its direction is mixed in. Equation \eqref{eq:diffmeans} is therefore also a recipe
for data collection: the two classes should differ in the concept and in nothing else. Pairs of prompts that
differ in a single word, like the clean and corrupted prompts of \cref{sec:interventions}, are built for this
purpose.

\begin{caution}
The hypothesis has known exceptions. \citet{engels2025not} found that GPT-2 and Mistral 7B represent the day of
the week and the month of the year on a circle inside a two-dimensional subspace. Such a feature is not a single
direction with a strength. There are also several inequivalent ways to make the word ``linear'' precise, which
\citet{park2024linear} disentangle. \Cref{def:lrh} should be read as a productive first approximation.
\end{caution}

\subsection{Superposition: more features than dimensions}
\label{sec:superposition}

\Cref{def:lrh} does not say how large $N$ is. If the directions had to be orthogonal, we would have $N\le d$, and
a stream of width $768$ could hold $768$ features. The world contains far more than $768$ things worth
representing. The \emph{superposition hypothesis} \citep{elhage2022toy} says that networks escape this limit by
using directions that are only \emph{almost} orthogonal, and that $N$ is much larger than $d$. A related picture
for word embeddings was proposed earlier by \citet{arora2018linear}. The hypothesis rests on a geometric fact:
in high dimension there is room for a number of almost orthogonal directions that grows exponentially with $d$.

\begin{proposition}[Many almost orthogonal directions]\label{prop:jl}
Let $0<\varepsilon<1$ and let $N\ge 2$ be an integer with $N\le\exp(\varepsilon^2d/4)$. Then there are unit
vectors $v_1,\dots,v_N\in\R^d$ with $|v_i^\top v_j|\le\varepsilon$ for all $i\ne j$.
\end{proposition}
\begin{proof}
We choose the vectors at random and show that they have the property with positive probability. Let all $dN$
coordinates of $v_1,\dots,v_N$ be independent, each equal to $+1/\sqrt d$ or $-1/\sqrt d$ with probability $1/2$.
Every $v_j$ then has norm one.

\textbf{One pair.} Fix $i\ne j$. The number $d\,v_i^\top v_j$ is a sum of $d$ terms, one for each coordinate, and
each term is the product of two independent uniform random signs. Such a product is again a uniform random
sign, and the $d$ terms are independent because the coordinates are. Hoeffding's inequality
\citep[Theorem 2]{hoeffding1963probability} states that a sum of $d$ independent random variables with values in
$[-1,1]$ and mean zero is at least $\varepsilon d$ with probability at most $\exp(-(\varepsilon d)^2/(2d))$. We
use it for the sum $d\,v_i^\top v_j$ and for its negative:
\begin{align}
\Pr\big(|v_i^\top v_j|>\varepsilon\big)
&\le \Pr\big(d\,v_i^\top v_j\ge\varepsilon d\big)+\Pr\big(-d\,v_i^\top v_j\ge\varepsilon d\big) && \by{union of two events}\notag\\
&\le 2\exp\!\big(-(\varepsilon d)^2/(2d)\big) && \by{Hoeffding, twice}\notag\\
&= 2\exp\!\big(-\varepsilon^2d/2\big). && \by{since $(\varepsilon d)^2/(2d)=\varepsilon^2d/2$} \label{eq:onepair}
\end{align}

\textbf{All pairs.} There are $N(N-1)/2$ pairs $i<j$. Hence
\begin{align*}
\Pr\big(\text{some pair has }|v_i^\top v_j|>\varepsilon\big)
&\le \frac{N(N-1)}{2}\cdot 2\exp\!\big(-\varepsilon^2d/2\big) && \by{union bound and \eqref{eq:onepair}}\\
&< N^2\exp\!\big(-\varepsilon^2d/2\big) && \by{since $N(N-1)<N^2$}\\
&\le \exp\!\big(\varepsilon^2d/2\big)\exp\!\big(-\varepsilon^2d/2\big) \;=\; 1. && \by{since $N\le\exp(\varepsilon^2d/4)$}
\end{align*}
The probability that some pair violates the bound is strictly smaller than one. So there is an outcome in
which no pair violates it, and this outcome is the required family of vectors.
\end{proof}

For $d=768$ and $\varepsilon=0.3$ the proposition allows $N\le e^{17.28}$, which is more than $3\times10^7$
directions, compared with $768$ exactly orthogonal ones. The constants in the bound are far from optimal. The
point is the exponential dependence on $d$. The mechanism is easy to remember. A sum of $d$ random signs has
typical size $\sqrt d$, so two random unit vectors have an inner product of typical size $1/\sqrt d$, which is
$0.036$ for $d=768$.

Almost orthogonal is not orthogonal, and the price is \emph{interference}. If we read feature $j$ from $x$ by
projecting on $v_j$, the other active features leak into the result.

\begin{lemma}[Interference is controlled by sparsity]\label{lem:interference}
Let $v_1,\dots,v_N$ be unit vectors with $|v_i^\top v_j|\le\varepsilon$ for $i\ne j$. Let $x=\sum_{i=1}^{N}f_iv_i$,
where at most $k$ of the numbers $f_1,\dots,f_N$ are nonzero and all of them lie in $[0,1]$. Then
$|v_j^\top x-f_j|\le k\varepsilon$ for every $j$.
\end{lemma}
\begin{proof}
We expand $v_j^\top x$ and bound the terms with $i\ne j$.
\begin{align*}
|v_j^\top x-f_j| &= \Big|\sum_{i=1}^{N}f_i\,v_j^\top v_i-f_j\Big| && \by{definition of $x$}\\
&= \Big|\sum_{i\ne j}f_i\,v_j^\top v_i\Big| && \by{since $v_j^\top v_j=1$}\\
&\le \sum_{i\ne j}f_i\,|v_j^\top v_i| && \by{triangle inequality, $f_i\ge0$}\\
&\le k\varepsilon. && \by{at most $k$ nonzero terms, each at most $\varepsilon$}\qedhere
\end{align*}
\end{proof}

The lemma shows what superposition needs: \emph{sparsity}. Packing many features into few dimensions is
harmless as long as only a few of them are active at the same time, because then $k\varepsilon$ is small. Most
features of language are of this kind. Almost no text is about the Golden Gate Bridge. A nonlinearity can
then remove the remaining interference. If feature $j$ is absent, the lemma gives $v_j^\top x\le k\varepsilon$,
so $\relu(v_j^\top x-k\varepsilon)=0$: a threshold reads an absent feature as exactly absent.

\begin{figure}[t]
\centering
\begin{tikzpicture}[>=Latex, font=\small]
% panel a
\begin{scope}
\draw[tutgrey!40] (-1.5,0) -- (1.5,0); \draw[tutgrey!40] (0,-1.5) -- (0,1.5);
\draw[->, line width=1.2pt, tutblue] (0,0) -- (1.25,0) node[below] {$v_1$};
\draw[->, line width=1.2pt, tutblue] (0,0) -- (0,1.25) node[right] {$v_2$};
\node[align=center] at (0,-2.15) {(a) two features,\\ two dimensions};
\end{scope}
% panel b
\begin{scope}[xshift=5.0cm]
\draw[tutgrey!40] (-1.5,0) -- (1.5,0); \draw[tutgrey!40] (0,-1.5) -- (0,1.5);
\foreach \k in {1,...,5} {\draw[->, line width=1.2pt, tutred] (0,0) -- ({1.25*cos(90+72*(\k-1))},{1.25*sin(90+72*(\k-1))});}
\node[tutred!80!black] at (0.35,1.38) {$v_1$};
\node[align=center] at (0,-2.15) {(b) five features,\\ two dimensions};
\end{scope}
% panel c
\begin{scope}[xshift=10.0cm]
\draw[tutgrey!40] (-1.5,0) -- (1.5,0);
\draw[->, line width=1.2pt, tutred] (0,0) -- (1.25,0) node[above] {$v_1$};
\draw[->, line width=1.2pt, tutred] (0,0) -- (-1.25,0) node[above] {$v_2$};
\fill[tutgrey] (0,0) circle (1.3pt);
\node[align=center] at (0,-2.15) {(c) two features,\\ one dimension};
\end{scope}
\end{tikzpicture}
\caption{Feature directions without superposition (a) and with superposition (b, c). In (b) neighbouring
directions of the pentagon have inner product $\cos 72^\circ\approx0.31$, and the others have
$\cos 144^\circ\approx-0.81$. In (c) the two directions are
antipodal, with inner product $-1$. Both arrangements appear as solutions of the toy model of
\citet{elhage2022toy}.}
\label{fig:superposition}
\end{figure}

\subsection{A toy model that can be solved by hand}
\label{sec:toy}

\citet{elhage2022toy} study superposition in the smallest network that can show it. We describe their model in
the notation of \cref{def:lrh}. They write $x$ for our $f$ and $h$ for our $x$, and their sparsity parameter is our $1-\delta$.

There are $N$ features with activations $f=(f_1,\dots,f_N)$. The features are independent. Each $f_j$ is zero
with probability $1-\delta$ and uniformly distributed on $[0,1]$ with probability $\delta$, so a small $\delta$
means sparse features. The model compresses $f$ into $d<N$ dimensions and tries to recover it:
\begin{equation}\label{eq:toy}
x \;=\; Wf\in\R^d,\qquad \hat f \;=\; \relu\big(W^\top x+b_{\mathrm{toy}}\big)\in\R^N,\qquad
\text{loss} \;=\; \E\sum_{j=1}^{N}\big(f_j-\hat f_j\big)^2 ,
\end{equation}
with a weight matrix $W\in\R^{d\times N}$ and a bias $b_{\mathrm{toy}}\in\R^N$. Column $j$ of $W$ is the direction
that the model assigns to feature $j$. The original model also weights each feature by an importance, which we
set to one.

We work out the case of $N=2$ features and $d=1$ dimension completely. The model has one number $x$ to store
two features. We compare two ways of using it.

\begin{itemize}
\item \textbf{Dedicated solution.} $W=(1,\ 0)$ and $b_{\mathrm{toy}}=(0,\ \delta/2)$. The model stores feature $1$
and ignores feature $2$, for which it always outputs the constant $\delta/2$.
\item \textbf{Antipodal solution.} $W=(1,\ -1)$ and $b_{\mathrm{toy}}=(0,\ 0)$. The model stores the difference
$x=f_1-f_2$, as in \cref{fig:superposition}(c). For example, the input $f=(0.6,\ 0.2)$ is stored as $x=0.4$ and read
back as $\hat f=(0.4,\ 0)$.
\end{itemize}

\begin{proposition}[Superposition wins when features are sparse]\label{prop:toy}
In the model \eqref{eq:toy} with $N=2$, $d=1$ and $0<\delta\le1$, the dedicated solution has loss
$\delta/3-\delta^2/4$ and the antipodal solution has loss $\delta^2/3$. The antipodal solution has the smaller
loss if and only if $\delta<4/7$.
\end{proposition}
\begin{proof}
We use that a variable which is uniform on $[0,1]$ has mean $1/2$ and second moment $1/3$. Hence
$\E f_2=\delta/2$ and $\E f_2^2=\delta/3$.

\textbf{Dedicated solution.} Here $x=f_1$, so $\hat f_1=\relu(f_1)=f_1$ and $\hat f_2=\relu(\delta/2)=\delta/2$.
\begin{align*}
\text{loss} &= \E\big(f_1-f_1\big)^2+\E\big(f_2-\delta/2\big)^2 && \by{inserting $\hat f_1$ and $\hat f_2$}\\
&= \E f_2^2-\delta\,\E f_2+\delta^2/4 && \by{expanding the square}\\
&= \delta/3-\delta^2/2+\delta^2/4 \;=\; \delta/3-\delta^2/4. && \by{since $\E f_2=\delta/2$, $\E f_2^2=\delta/3$}
\end{align*}

\textbf{Antipodal solution.} Here $x=f_1-f_2$, so $\hat f_1=\relu(f_1-f_2)$ and $\hat f_2=\relu(f_2-f_1)$. If
$f_2=0$, then $\hat f_1=f_1$ and $\hat f_2=\relu(-f_1)=0=f_2$, and the error is zero. The same holds if $f_1=0$.
The only inputs with an error are those with both features active, which has probability $\delta^2$. Let both
be active and let $f_1\ge f_2$. Then $\hat f_1=f_1-f_2$ and $\hat f_2=0$, so both errors equal $f_2$ and the
squared error is $2f_2^2=2\min(f_1,f_2)^2$. For $f_2\ge f_1$ we obtain the same expression by symmetry. Given that
both features are active, $f_1$ and $f_2$ are independent and uniform on $[0,1]$, so
$\Pr(\min(f_1,f_2)>u)=(1-u)^2$ and the minimum has density $2(1-u)$ on $[0,1]$.
\begin{align*}
\text{loss} &= \delta^2\cdot 2\int_0^1u^2\cdot 2(1-u)\,du && \by{error only if both active}\\
&= \delta^2\cdot 4\Big(\frac13-\frac14\Big) \;=\; \frac{\delta^2}{3}. && \by{since $\textstyle\int_0^1u^2du=\frac13$, $\int_0^1u^3du=\frac14$}
\end{align*}

\textbf{Comparison.} For $\delta>0$,
\begin{align*}
\frac{\delta^2}{3}<\frac{\delta}{3}-\frac{\delta^2}{4}
\;&\Longleftrightarrow\; \frac{\delta}{3}+\frac{\delta}{4}<\frac13 && \by{dividing by $\delta$, rearranging}\\
&\Longleftrightarrow\; \delta<\frac47. && \by{since $\tfrac13+\tfrac14=\tfrac{7}{12}$}\qedhere
\end{align*}
\end{proof}

\begin{intuition}
The antipodal solution is wrong only when both features are present at once, and in that case it is badly
wrong. The dedicated solution is always exactly right about one feature and always ignorant about the other.
If features are rare, the collisions are rarer still ($\delta^2$ against $\delta$), and it pays to store both.
Superposition is a bet on sparsity.
\end{intuition}

Neither of the two solutions needs to be the global minimum of the loss, and the proposition only compares the
two strategies. \citet{elhage2022toy} minimise the loss numerically for larger $N$ and $d$ and find the same
transition. When the features are dense, the model stores the $d$ most important ones in orthogonal directions
and drops the rest. As the features become sparser, it stores more and more of them, in regular arrangements
such as antipodal pairs and the pentagon of \cref{fig:superposition}(b).

\subsection{Consequences for interpretability}

If superposition is how networks store features, three conclusions follow.

\begin{itemize}
\item \textbf{Polysemantic neurons are expected.} With $N$ much larger than the number of neurons, the feature
directions cannot all be aligned with the neuron basis. Each neuron then has a nonzero component along many
feature directions and responds to all of the corresponding features.
\item \textbf{The right unit of analysis has to be found.} The features are directions that no part of the
architecture labels for us. \citet{gurnee2023finding} show both situations in real models: some features sit in
single neurons, and others are spread over many.
\item \textbf{Finding them is a sparse recovery problem.} We observe many activation vectors $x$. We look for a
set of directions, more numerous than the dimension, such that each $x$ is a combination of a few of them. This
problem is known as sparse dictionary learning, and \cref{sec:sae} describes the method that is currently used
for it.
\end{itemize}

\Cref{lem:interference} concerns the \emph{storage} of features. Whether and how networks can also
\emph{compute} with features in superposition is a separate question with its own theory
\citep{hanni2024mathematical,adler2024complexity}.

\begin{takeaways}
The working hypothesis is that features are directions, and that their activations are lengths along these
directions. High-dimensional spaces contain exponentially many almost orthogonal directions, so a network can
store far more features than it has dimensions. The interference between them is small as long as few features
are active at once. Neurons are then mixtures of many features, and the features have to be recovered by sparse
dictionary learning.
\end{takeaways}

\begin{exercise}\label{ex:triangle}
Three features are stored in $d=2$ dimensions along unit vectors at mutual angles of $120^\circ$. Show that
$v_i^\top v_j=-1/2$ for $i\ne j$. Show that the read-out $\hat f_j=\relu(v_j^\top x)$ recovers $f$ exactly whenever
at most one feature is active.
\end{exercise}
\begin{exercise}\label{ex:bound}
Evaluate the bound of \cref{prop:jl} for $d=4096$ with $\varepsilon=0.1$ and with $\varepsilon=0.2$. For which
$k$ does \cref{lem:interference} then guarantee an interference below $1/2$?
\end{exercise}
\begin{exercise}\label{ex:nobias}
In \cref{prop:toy}, replace the bias of the dedicated solution by $b_{\mathrm{toy}}=(0,0)$. Compute the loss and
show that the antipodal solution is now better for every $\delta<1$. What does the bias $\delta/2$ achieve?
\end{exercise}
